\documentclass[11pt,a4paper]{article}
\usepackage[utf8]{inputenc}
\usepackage[indonesian]{babel}
\usepackage{amsmath,amsfonts,amssymb}
\usepackage{geometry}
\usepackage[most]{tcolorbox}
\usepackage{xcolor}
\usepackage{fancyhdr}
\usepackage{booktabs}
\usepackage{array}
\usepackage{enumitem}
\usepackage{hyperref}
\usepackage{parskip}
\usepackage{adjustbox}

\geometry{margin=2.5cm, headheight=14pt, top=3cm}

\definecolor{primary}{HTML}{1E3A5F}
\definecolor{accent}{HTML}{EBF5FB}
\definecolor{gold}{HTML}{D4AC0D}

\tcbset{
  formula_box/.style={
    colback=accent, colframe=primary,
    fonttitle=\bfseries\color{white},
    coltitle=white, attach boxed title to top left={yshift=-2mm,xshift=6mm},
    boxed title style={colback=primary, rounded corners},
    boxrule=0.8pt, arc=6pt,
    left=8pt, right=8pt, top=16pt, bottom=8pt,
    after skip=14pt
  }
}

\pagestyle{fancy}
\fancyhf{}
\lhead{\small\bfseries Ammar Hanafi (2206051582)}
\rhead{\small\textit{Analisis STATIS -- Jawa Barat, Banten, DKI Jakarta}}
\cfoot{\small\thepage}
\renewcommand{\headrulewidth}{0.4pt}

\setlist[enumerate]{leftmargin=*, itemsep=2pt, topsep=4pt}
\setlist[itemize]{leftmargin=*, itemsep=2pt, topsep=4pt}

\begin{document}

\begin{titlepage}
  \centering
  \vspace*{3cm}
  {\Huge\bfseries Penjabaran Langkah demi Langkah\\[0.5em]
  Seluruh Formula Matematis\\[1em]}
  {\Large\textit{Analisis STATIS Panel Sosio-Ekonomi\\
  41 Kabupaten/Kota, 5 Tahun (2020--2024)}}\\[3em]
  \rule{0.6\textwidth}{0.5pt}\\[1.5em]
  {\large Ammar Hanafi\\NPM: 2206051582\\[0.5em]
  Program Studi Statistika\\
  Fakultas Matematika dan Ilmu Pengetahuan Alam\\
  Universitas Indonesia\\[1.5em]
  \textbf{Terakhir Diperbarui: \today}}
  \vspace*{3cm}
\end{titlepage}

\tableofcontents
\newpage

\section*{Pendahuluan} \label{sec:method}
Dokumen ini menjabarkan \textbf{setiap formula matematis} yang digunakan dalam penelitian skripsi secara lengkap dengan:
\begin{enumerate}
  \item \textbf{Formula Asal} --- Notasi simbolik baku
  \item \textbf{Keterangan Variabel} --- Definisi setiap simbol
  \item \textbf{Substitusi Langkah demi Langkah} --- Nilai \emph{nyata} dari \texttt{data2.xlsx} dimasukkan
  \item \textbf{Hasil Akhir} --- Nilai komputasi sesungguhnya
\end{enumerate}

\noindent Data yang digunakan: \textbf{41 Kabupaten/Kota} (Provinsi Banten, DKI Jakarta, Jawa Barat), \textbf{12 variabel sosio-ekonomi}, periode \textbf{2020--2024} ($T = 5$ tahun).

\subsection*{Daftar Singkatan dan Variabel}
\begin{adjustbox}{max width=\textwidth}
\begin{tabular}{lll}
\toprule
\textbf{Singkatan} & \textbf{Kepanjangan} & \textbf{Satuan} \\
\midrule
IPM   & Indeks Pembangunan Manusia & Indeks \\
IPG   & Indeks Pembangunan Gender  & Indeks \\
PPM   & Persentase Penduduk Miskin & Persen (\%) \\
PDRB\_ADHK & Produk Domestik Regional Bruto Atas Dasar Harga Konstan & Juta Rupiah \\
RLS   & Rata-rata Lama Sekolah    & Tahun \\
AHH\_L & Angka Harapan Hidup Laki-Laki & Tahun \\
AHH\_P & Angka Harapan Hidup Perempuan & Tahun \\
GK    & Garis Kemiskinan          & Rupiah/Kapita/Bulan \\
PKKP  & Pengeluaran per Kapita yang Disesuaikan & Ribu Rupiah/Orang/Tahun \\
TPT   & Tingkat Pengangguran Terbuka & Persen (\%) \\
TPAK  & Tingkat Partisipasi Angkatan Kerja & Persen (\%) \\
GR    & Gini Ratio (Rasio Ketimpangan) & Indeks [0,1] \\
\bottomrule
\end{tabular}
\end{adjustbox}

\subsection*{Notasi Umum}
\begin{adjustbox}{max width=\textwidth}
\begin{tabular}{ll}
\toprule
\textbf{Notasi} & \textbf{Arti} \\
\midrule
$n$ & Jumlah wilayah observasi ($n = 41$) \\
$p$ & Jumlah variabel ($p = 12$) \\
$T$ & Jumlah tahun pengamatan ($T = 5$: 2020--2024) \\
$i$ & Indeks wilayah ($i = 1, 2, \ldots, n$) \\
$j$ & Indeks variabel ($j = 1, 2, \ldots, p$) \\
$t$ & Indeks tahun ($t = 1, 2, \ldots, T$) \\
$\mathbf{X}_t$ & Matriks data mentah ($n \times p$) tahun ke-$t$ \\
$\mathbf{Z}_t$ & Matriks data terstandarisasi ($n \times p$) tahun ke-$t$ \\
$\mathbf{W}_t$ & Matriks kedekatan antar-wilayah ($n \times n$) tahun ke-$t$ \\
$\mathbf{W}_c$ & Matriks kompromi STATIS ($n \times n$) \\
$\mathbf{F}_c$ & Matriks koordinat skor kompromi ($n \times 2$) \\
KMO & Kaiser-Meyer-Olkin Measure of Sampling Adequacy \\
MSA & Measure of Sampling Adequacy (per variabel) \\
RV & Koefisien kemiripan antar-matriks (Escoufier) \\
PCA & Principal Component Analysis (Analisis Komponen Utama) \\
SVD & Singular Value Decomposition (Dekomposisi Nilai Singular) \\
SSE & Sum of Squared Errors (Jumlah Kuadrat Galat) \\
STATIS & Structuration des Tableaux \`A Trois Indices de la Statistique \\
\bottomrule
\end{tabular}
\end{adjustbox}

\newpage

\section{Analisis Statistika Deskriptif}
\begin{tcolorbox}[formula_box, title={1. Standardisasi Data (Z-Score)}]
\textbf{Formula:}
\[ z_{ij} = \frac{x_{ij} - \bar{x}_j}{s_j} \]

\textbf{Keterangan Variabel:}
\begin{itemize}
  \item $x_{ij}$: Nilai mentah variabel $j$ pada wilayah ke-$i$
  \item $\bar{x}_j$: Rata-rata variabel $j$ dari seluruh $n=41$ wilayah
  \item $s_j$: Simpangan baku variabel $j$
\end{itemize}

\textbf{Substitusi (Langkah demi Langkah):}
\begin{enumerate}
  \item Ambil contoh: $i = 1$ (Kep.\ Seribu), $j = $ IPM, tahun 2020
  \item $x_{1,IPM} = 71.63$
  \item $\bar{x}_{IPM} = \frac{\sum_{i=1}^{41} x_{i,IPM}}{41} = \frac{2988.78}{41} = 72.8971$
  \item $s_{IPM} = 5.7902$
  \item $z_{1,IPM} = \frac{71.63 - 72.8971}{5.7902} = \frac{-1.2671}{5.7902}$
\end{enumerate}

\textbf{Hasil Akhir:}
\[ z_{1,IPM} = -0.2188 \]
\textit{(Nilai negatif: IPM Kep.\ Seribu di bawah rata-rata)}
\end{tcolorbox}

\begin{tcolorbox}[formula_box, title={2. Rata-Rata (Mean)}]
\textbf{Formula:}
\[ \bar{x}_j = \frac{1}{n} \sum_{i=1}^{n} x_{ij} \]

\textbf{Keterangan Variabel:}
\begin{itemize}
  \item $\bar{x}_j$: Rata-rata variabel $j$
  \item $n$: Jumlah wilayah pengamatan
  \item $x_{ij}$: Nilai variabel $j$ untuk wilayah ke-$i$
\end{itemize}

\textbf{Substitusi (Langkah demi Langkah):}
\begin{enumerate}
  \item Variabel: IPM, Tahun 2020, $n = 41$ wilayah
  \item $\bar{x}_{IPM} = \frac{1}{41} \sum_{i=1}^{41} x_{i,IPM}$
  \item $= \frac{71.63 + 84.72 + 82.66 + \cdots}{41}$
  \item $= \frac{2988.78}{41}$
\end{enumerate}

\textbf{Hasil Akhir:}
\[ \bar{x}_{IPM} = 72.8971 \]
\end{tcolorbox}

\begin{tcolorbox}[formula_box, title={3. Median (Kuartil Kedua, $Q_2$)}]
\textbf{Formula:}
\[ Q_2 = \begin{cases} x_{\left(\frac{n+1}{2}\right)} & n \text{ ganjil} \\ \frac{x_{\left(\frac{n}{2}\right)} + x_{\left(\frac{n}{2}+1\right)}}{2} & n \text{ genap} \end{cases} \]

\textbf{Keterangan Variabel:}
\begin{itemize}
  \item $n = 41$ (ganjil), sehingga $Q_2 = x_{(21)}$
  \item $x_{(k)}$: Nilai ke-$k$ setelah data diurutkan naik
\end{itemize}

\textbf{Substitusi (Langkah demi Langkah):}
\begin{enumerate}
  \item Data IPM 2020 diurutkan: $\{63.91, 65.00, \ldots, 84.72\}$
  \item $n = 41$ (ganjil) $\Rightarrow$ posisi median $= \frac{41+1}{2} = 21$
  \item $Q_2 = x_{(21)} = 71.7000$
\end{enumerate}

\textbf{Hasil Akhir:}
\[ Q_2 = 71.7000 \]
\end{tcolorbox}

\begin{tcolorbox}[formula_box, title={4. Simpangan Baku ($s$)}]
\textbf{Formula:}
\[ s_j = \sqrt{\frac{1}{n-1} \sum_{i=1}^{n} (x_{ij} - \bar{x}_j)^2} \]

\textbf{Keterangan Variabel:}
\begin{itemize}
  \item $s_j$: Simpangan baku sampel variabel $j$
  \item $n-1$: Derajat kebebasan (koreksi Bessel)
\end{itemize}

\textbf{Substitusi (Langkah demi Langkah):}
\begin{enumerate}
  \item Variabel IPM, 2020: $\bar{x}_{IPM} = 72.8971$
  \item $\sum_{i=1}^{41} (x_{i,IPM} - \bar{x})^2 = 1341.0352$
  \item $s^2 = \frac{1341.0352}{41-1} = \frac{1341.0352}{40} = 33.5259$
  \item $s = \sqrt{33.5259}$
\end{enumerate}

\textbf{Hasil Akhir:}
\[ s_{IPM} = 5.7902 \]
\end{tcolorbox}

\begin{tcolorbox}[formula_box, title={5. Nilai Minimum dan Maksimum}]
\textbf{Formula:}
\[ x_{\min} = \min(x_1,\ldots,x_n),\quad x_{\max} = \max(x_1,\ldots,x_n) \]

\textbf{Keterangan Variabel:}
\begin{itemize}
  \item $x_{\min}$ dan $x_{\max}$: Nilai ekstrem dalam distribusi variabel
\end{itemize}

\textbf{Substitusi (Langkah demi Langkah):}
\begin{enumerate}
  \item Variabel IPM, 2020 (41 wilayah)
  \item Survey seluruh nilai: IPM terendah $= 63.91$, tertinggi $= 84.72$
\end{enumerate}

\textbf{Hasil Akhir:}
\[ x_{\min,IPM} = 63.91 \quad x_{\max,IPM} = 84.72 \]
\textit{(Lebak hingga Kota Jakarta Selatan)}
\end{tcolorbox}

\begin{tcolorbox}[formula_box, title={6. Jangkauan Interkuartil (IQR)}]
\textbf{Formula:}
\[ IQR = Q_3 - Q_1 \]

\textbf{Keterangan Variabel:}
\begin{itemize}
  \item $Q_1$: Kuartil pertama (persentil ke-25)
  \item $Q_3$: Kuartil ketiga (persentil ke-75)
\end{itemize}

\textbf{Substitusi (Langkah demi Langkah):}
\begin{enumerate}
  \item Variabel IPM, 2020
  \item $Q_1 = 68.0800$ (persentil ke-25)
  \item $Q_3 = 77.8300$ (persentil ke-75)
  \item $IQR = 77.8300 - 68.0800$
\end{enumerate}

\textbf{Hasil Akhir:}
\[ IQR_{IPM} = 9.7500 \]
\end{tcolorbox}

\begin{tcolorbox}[formula_box, title={7. Koefisien Variasi / \textit{Coefficient of Variation} (CV)}]
\textbf{Formula:}
\[ CV = \frac{s}{\bar{x}} \times 100\% \]

\textbf{Keterangan Variabel:}
\begin{itemize}
  \item $CV$ (\textit{Coefficient of Variation}): Ukuran variabilitas relatif, bebas satuan
  \item $s$: Simpangan baku; $\bar{x}$: Rata-rata
  \item Klasifikasi (Sumber: Hasan, 2001): $CV \leq 10\%$ = rendah; $10\% < CV \leq 20\%$ = sedang; $CV > 20\%$ = tinggi
\end{itemize}

\textbf{Substitusi (Langkah demi Langkah):}
\begin{enumerate}
  \item Variabel IPM, 2020: $s = 5.7902$, $\bar{x} = 72.8971$
  \item $CV = \frac{5.7902}{72.8971} \times 100\%$
  \item $CV = 0.0794 \times 100\%$
\end{enumerate}

\textbf{Hasil Akhir:}
\[ CV_{IPM} = 7.94\% \quad \text{(rendah ($\leq 10\%$))} \]
\end{tcolorbox}

\begin{tcolorbox}[formula_box, title={8. Koefisien Kemiringan / \textit{Skewness} ($g_1$)}]
\textbf{Formula:}
\[ g_1 = \frac{\frac{1}{n}\sum_{i=1}^{n}(x_i-\bar{x})^3}{s^3} \]

\textbf{Keterangan Variabel:}
\begin{itemize}
  \item $g_1$ (huruf $g$ dari \textit{gamma}, notasi standar Fisher): Koefisien kemiringan
  \item Klasifikasi (Bulmer, 1979): $|g_1| < 0.5$: simetri; $0.5 \leq |g_1| < 1$: menceng moderat; $|g_1| \geq 1$: sangat menceng
  \item $g_1 > 0$: Distribusi menceng kanan (\textit{right-skewed})
  \item $g_1 < 0$: Distribusi menceng kiri (\textit{left-skewed})
\end{itemize}

\textbf{Substitusi (Langkah demi Langkah):}
\begin{enumerate}
  \item $\frac{1}{n}\sum(x_i-\bar{x})^3 = 81.9889$
  \item $s^3 = 5.7902^3 = 194.1200$
  \item $g_1 = \frac{81.9889}{194.1200}$
\end{enumerate}

\textbf{Hasil Akhir:}
\[ g_{1,IPM} = 0.4224 \quad \text{(simetri)} \]
\end{tcolorbox}

\begin{tcolorbox}[formula_box, title={9. Koefisien Keruncingan / \textit{Excess Kurtosis} ($g_2$)}]
\textbf{Formula:}
\[ g_2 = \frac{\frac{1}{n}\sum_{i=1}^{n}(x_i-\bar{x})^4}{s^4} - 3 \]

\textbf{Keterangan Variabel:}
\begin{itemize}
  \item $g_2$ (huruf $g$ dari \textit{gamma}, notasi standar Fisher): Koefisien keruncingan
  \item Pengurangan $-3$ disebut \textit{excess kurtosis} (Westfall, 2014): kurtosis distribusi normal $= 3$, sehingga dikurangi 3 agar baseline $= 0$
  \item $g_2 > 0$: Leptokurtik (lebih runcing, ekor lebih tebal dari normal)
  \item $g_2 < 0$: Platikurtik (lebih datar, ekor lebih tipis dari normal)
\end{itemize}

\textbf{Substitusi (Langkah demi Langkah):}
\begin{enumerate}
  \item $\frac{1}{n}\sum(x_i-\bar{x})^4 = 2166.7401$
  \item $s^4 = 5.7902^4 = 1123.9847$
  \item $g_2 = \frac{2166.7401}{1123.9847} - 3 = 1.9277 - 3$
\end{enumerate}

\textbf{Hasil Akhir:}
\[ g_{2,IPM} = -1.0723 \quad \text{(platikurtik)} \]
\end{tcolorbox}


\section{Uji Kelayakan Data}
\begin{tcolorbox}[formula_box, title={10. Kaiser-Meyer-Olkin (KMO)}]
\textbf{Formula:}
\[ KMO = \frac{\sum_{j \neq k} r_{jk}^2}{\sum_{j \neq k} r_{jk}^2 + \sum_{j \neq k} r_{jk.p}^2} \]

\textbf{Keterangan Variabel:}
\begin{itemize}
  \item $r_{jk}$: Korelasi Pearson antara variabel $j$ dan $k$
  \item $r_{jk.p}$: Korelasi Parsial antara $j$ dan $k$ (setelah dikontrol variabel lain)
  \item Klasifikasi KMO (Kaiser, 1974): $<0.5$ = \textit{Unacceptable}; $0.5$--$0.6$ = \textit{Miserable}; $0.6$--$0.7$ = \textit{Mediocre}; $0.7$--$0.8$ = \textit{Middling}; $0.8$--$0.9$ = \textit{Meritorious}; $>0.9$ = \textit{Marvelous}
\end{itemize}

\textbf{Substitusi (Langkah demi Langkah):}
\begin{enumerate}
  \item Dari matriks korelasi $\mathbf{R}$ ($p=12$ variabel, $n=41$ observasi)
  \item $\sum_{j\neq k} r_{jk}^2 = 32.2239$
  \item $\sum_{j\neq k} r_{jk.p}^2 = 8.7302$
  \item $KMO = \frac{32.2239}{32.2239 + 8.7302} = \frac{32.2239}{40.9542}$
\end{enumerate}

\textbf{Hasil Akhir:}
\[ KMO = 0.7868 \quad \text{(\textit{Middling}, Kaiser 1974)} \]
\end{tcolorbox}

\begin{tcolorbox}[formula_box, title={11. Derajat Kebebasan Uji Bartlett ($df$)}]
\textbf{Formula:}
\[ df = \frac{p(p-1)}{2} \]

\textbf{Keterangan Variabel:}
\begin{itemize}
  \item $p$: Jumlah variabel; $df$: Jumlah pasangan unik variabel
\end{itemize}

\textbf{Substitusi (Langkah demi Langkah):}
\begin{enumerate}
  \item $p = 12$ variabel
  \item $df = \frac{12 \times (12-1)}{2} = \frac{132}{2}$
\end{enumerate}

\textbf{Hasil Akhir:}
\[ df = 66 \]
\end{tcolorbox}

\begin{tcolorbox}[formula_box, title={12. Uji Bartlett's Test of Sphericity ($\chi^2$)}]
\textbf{Formula:}
\[ \chi^2 = -\left[n - 1 - \frac{2p+5}{6}\right] \ln|\mathbf{R}| \]

\textbf{Keterangan Variabel:}
\begin{itemize}
  \item $\mathbf{R}$: Matriks korelasi $p \times p$ (dihitung dari data terstandarisasi)
  \item $|\mathbf{R}|$: Determinan matriks korelasi --- dihitung oleh \texttt{numpy.linalg.det(R)} dari matriks $12 \times 12$
  \item Jika $|\mathbf{R}| = 1$, maka variabel tidak berkorelasi. Semakin kecil $|\mathbf{R}|$, semakin tinggi korelasi
  \item $H_0$: $\mathbf{R} = \mathbf{I}$ (Sumber: Bartlett, 1950)
\end{itemize}

\textbf{Substitusi (Langkah demi Langkah):}
\begin{enumerate}
  \item $n = 41$, $p = 12$
  \item Faktor koreksi: $n-1-\frac{2p+5}{6} = 41-1-\frac{2(12)+5}{6} = 35.1667$
  \item $|\mathbf{R}| = 1.78e-08$ (dihitung dari matriks korelasi $12\times12$ data 2020)
  \item $\ln|\mathbf{R}| = \ln(1.78e-08) = -17.8424$
  \item $\chi^2 = -(35.1667) \times (-17.8424)$
\end{enumerate}

\textbf{Hasil Akhir:}
\[ \chi^2 = 627.4563, \quad df = 66, \quad p = 0 \]
\textit{(Tolak $H_0$: Terdapat korelasi signifikan, data sesuai untuk STATIS --- Bartlett, 1950)}
\end{tcolorbox}


\section{Korelasi dan Analisis Komponen Utama (PCA)}
\begin{tcolorbox}[formula_box, title={21. Korelasi Pearson ($r_{xy}$)}]
\textbf{Formula:}
\[ r_{xy} = \frac{\sum(x_i-\bar{x})(y_i-\bar{y})}{\sqrt{\sum(x_i-\bar{x})^2}\cdot\sqrt{\sum(y_i-\bar{y})^2}} \]

\textbf{Keterangan Variabel:}
\begin{itemize}
  \item $r_{xy}\in[-1,1]$: Ukuran hubungan linier
  \item $|r|>0.7$: Korelasi kuat
\end{itemize}

\textbf{Substitusi (Langkah demi Langkah):}
\begin{enumerate}
  \item IPM vs RLS, 2020: $\bar{x}=72.8971$, $\bar{y}=8.9244$
  \item Pembilang: $368.3503$
  \item Penyebut: $382.6320$
\end{enumerate}

\textbf{Hasil Akhir:}
\[ r_{IPM,RLS} = 0.9627 \]
\end{tcolorbox}

\begin{tcolorbox}[formula_box, title={22. Determinan Matriks ($|\mathbf{A}|$)}]
\textbf{Formula:}
\[ |\mathbf{A}| = \sum_{\sigma \in S_p} \text{sgn}(\sigma) \prod_{i=1}^{p} a_{i,\sigma(i)} \]

\textbf{Keterangan Variabel:}
\begin{itemize}
  \item $|\mathbf{A}|$: Determinan dari matriks bujur sangkar $\mathbf{A}$
  \item $S_p$: Himpunan seluruh permutasi indeks $\{1, \dots, p\}$
  \item $\text{sgn}(\sigma)$: Tanda (paritas) permutasi ($+1$ genap, $-1$ ganjil)
  \item $a_{i,\sigma(i)}$: Elemen matriks pada baris ke-$i$ dan kolom ke-$\sigma(i)$
  \item Digunakan menguji multikolinieritas. Semakin mendekati 0, korelasi makin kuat.
\end{itemize}

\textbf{Substitusi (Langkah demi Langkah):}
\begin{enumerate}
  \item Matriks Korelasi $\mathbf{R}$ berukuran $12\times12$ (data 2020)
  \item Dihitung produk permutasi elemen $\mathbf{R}$ lewat dekomposisi LU
\end{enumerate}

\textbf{Hasil Akhir:}
\[ |\mathbf{R}| = 1.78e-08 \]
\end{tcolorbox}

\begin{tcolorbox}[formula_box, title={23. PCA ($\mathbf{F}=\mathbf{Z}\mathbf{V}$)}]
\textbf{Formula:}
\[ \mathbf{F} = \mathbf{Z}\mathbf{V} \]

\textbf{Keterangan Variabel:}
\begin{itemize}
  \item $\mathbf{Z}$: Matriks terstandarisasi ($n\times p$)
  \item $\mathbf{V}$: Matriks eigenvector dari matriks korelasi $\mathbf{R}$
  \item $\mathbf{F}$: Skor komponen utama ($n\times k$)
\end{itemize}

\textbf{Substitusi (Langkah demi Langkah):}
\begin{enumerate}
  \item Dari dekomposisi matriks korelasi $\mathbf{R}$ (data 2020), diperoleh eigenvalue $\lambda_1=6.1267$, $\lambda_2=1.7915$
  \item Varians yang dijelaskan PC1 $= \lambda_1 / p \times 100\% = 51.1\%$
  \item Varians yang dijelaskan PC2 $= \lambda_2 / p \times 100\% = 14.9\%$
  \item Misal untuk Kep. Seribu (baris 1): $\mathbf{F}_{[1,:]} = \mathbf{Z}_{[1,:]} \cdot \mathbf{V}$
\end{enumerate}

\textbf{Hasil Akhir:}
\[ \begin{pmatrix} F_1 \\ F_2 \end{pmatrix}_{\text{Kep.Seribu}} = \begin{pmatrix} -1.1919 \\ -1.4198 \end{pmatrix} \quad \text{(Kumulatif: 66.0\%)} \]
\end{tcolorbox}

\begin{tcolorbox}[formula_box, title={24. Dekomposisi Eigen}]
\textbf{Formula:}
\[ \mathbf{R}\mathbf{v}=\lambda\mathbf{v} \]

\textbf{Keterangan Variabel:}
\begin{itemize}
  \item $\mathbf{R}$: Matriks korelasi data
  \item $\lambda$: Eigenvalue (besaran varians di sepanjang sumbu utama)
  \item $\mathbf{v}$: Eigenvector (arah komponen utama)
\end{itemize}

\textbf{Substitusi (Langkah demi Langkah):}
\begin{enumerate}
  \item Selesaikan persamaan karakteristik: $|\mathbf{R} - \lambda\mathbf{I}| = 0$ ($p=12$)
  \item Diperoleh 12 eigenvalue (diurutkan mengecil):
  \item $\lambda_1=6.1267, \lambda_2=1.7915, \lambda_3=1.4876$
  \item Total varians = $\sum\lambda_j = \text{tr}(\mathbf{R}) = 12$
\end{enumerate}

\textbf{Hasil Akhir:}
\[ \lambda = [6.127, 1.791, 1.488, 0.696, 0.538, \ldots] \]
\end{tcolorbox}

\begin{tcolorbox}[formula_box, title={25. Rotasi Baseline}]
\textbf{Formula:}
\[ \theta=-\operatorname{atan2}(v_2,v_1),\; \mathbf{Q}=\begin{pmatrix}\cos\theta&-\sin\theta\\\sin\theta&\cos\theta\end{pmatrix},\; \mathbf{F}^*=\mathbf{F}\mathbf{Q} \]

\textbf{Keterangan Variabel:}
\begin{itemize}
  \item $\theta$: Sudut rotasi untuk membuat Kep. Seribu sejajar sumbu X
  \item $\mathbf{Q}$: Matriks rotasi ortogonal $2\times2$
\end{itemize}

\textbf{Substitusi (Langkah demi Langkah):}
\begin{enumerate}
  \item Koordinat Kompromi Kep. Seribu: $v_1=-0.4150$, $v_2=0.4857$
  \item $\theta = -\text{atan2}(0.4857, -0.4150) = -130.51^\circ$
  \item $\mathbf{Q} = \begin{pmatrix} -0.6496 & 0.7602 \\ -0.7602 & -0.6496 \end{pmatrix}$
  \item $\mathbf{F}^* = [-0.4150, 0.4857] \times \mathbf{Q}$
\end{enumerate}

\textbf{Hasil Akhir:}
\[ \mathbf{F}^*_{\text{Kep.Seribu}} = [-0.0996, -0.6310] \]
\end{tcolorbox}


\section{Tahapan Metode STATIS}
\begin{tcolorbox}[formula_box, title={13. Intra-struktur: Matriks Kemiripan Tahunan ($\mathbf{W}_t$)}]
\textbf{Formula:}
\[ \mathbf{W}_t = \mathbf{Z}_t \mathbf{Z}_t^T \]

\textbf{Keterangan Variabel:}
\begin{itemize}
  \item $\mathbf{Z}_t$: Matriks data z-score ($n \times p$) tahun ke-$t$
  \item $\mathbf{W}_t$: Matriks simetri ($n \times n$) yang merepresentasikan kedekatan antar wilayah di tahun $t$
  \item Setiap elemen $w_{ij,t}$ = inner product baris $i$ dan $j$ dari $\mathbf{Z}_t$
\end{itemize}

\textbf{Substitusi (Langkah demi Langkah):}
\begin{enumerate}
  \item Tahun 2020: $\mathbf{Z}_{2020}$ berukuran $41 \times 12$
  \item $\mathbf{W}_{2020} = \mathbf{Z}_{2020} \cdot \mathbf{Z}_{2020}^T$
  \item Dimensi hasil: $(41 \times 12)(12 \times 41) = 41 \times 41$
  \item Elemen pojok kiri atas $\mathbf{W}_{2020}[1,1] = 1.3923$ (norma kuadrat Kep.\ Seribu)
\end{enumerate}

\textbf{Hasil Akhir:}
\[ \mathbf{W}_{2020} = 41 \times 41 \text{ matriks (nilai diagonal}\ [0,0] = 1.3923) \]
\end{tcolorbox}

\begin{tcolorbox}[formula_box, title={14. Koefisien RV (Inter-Struktur)}]
\textbf{Formula:}
\[ RV(\mathbf{W}_i,\mathbf{W}_j) = \frac{\text{tr}(\mathbf{W}_i\mathbf{W}_j)}{\sqrt{\text{tr}(\mathbf{W}_i^2)\,\text{tr}(\mathbf{W}_j^2)}} \]

\textbf{Keterangan Variabel:}
\begin{itemize}
  \item $\text{tr}(\cdot)$: Trace (jumlah elemen diagonal)
  \item $RV \in [0,1]$: Ukuran kesamaan struktur dua matriks
  \item $RV$ tinggi $\Rightarrow$ struktur antar-tahun stabil
\end{itemize}

\textbf{Substitusi (Langkah demi Langkah):}
\begin{enumerate}
  \item Contoh: Tahun 2020 vs 2021
  \item $\text{tr}(\mathbf{W}_{2020} \cdot \mathbf{W}_{2021}) = 471.5136$
  \item $\text{tr}(\mathbf{W}_{2020}^2) = 516.2532$
  \item $\text{tr}(\mathbf{W}_{2021}^2) = 473.9093$
  \item $RV = \frac{471.5136}{\sqrt{516.2532 \times 473.9093}} = \frac{471.5136}{494.6283}$
\end{enumerate}

\textbf{Hasil Akhir:}
\[ RV(\mathbf{W}_{2020}, \mathbf{W}_{2021}) = 0.9533 \quad \text{(Sangat stabil)} \]
\end{tcolorbox}

\begin{tcolorbox}[formula_box, title={15. Bobot Optimal ($\alpha_t$)}]
\textbf{Formula:}
\[ \alpha_t = \frac{u_{1t}}{\sqrt{\sum_{t=1}^{T} u_{1t}^2}} \]

\textbf{Keterangan Variabel:}
\begin{itemize}
  \item $\mathbf{u}_1$: Eigenvector pertama dari matriks inter-struktur $\mathbf{R}_{RV}$
  \item $u_{1t}$: Komponen ke-$t$ dari eigenvector $\mathbf{u}_1$
  \item $\alpha_t$: Bobot (kontribusi) tahun ke-$t$ dalam kompromi
\end{itemize}

\textbf{Substitusi (Langkah demi Langkah):}
\begin{enumerate}
  \item Eigenvektor pertama $\mathbf{u}_1$ dari $\mathbf{R}_{RV}$ ($T=5$ tahun):
  \item $\mathbf{u}_1 = [0.4480, 0.4451, 0.4432, 0.4508, 0.4489]$
  \item Penyebut: $\sqrt{(0.4480)^2 + (0.4451)^2 + (0.4432)^2 + (0.4508)^2 + (0.4489)^2} = \sqrt{1.0000}$
  \item Masing-masing $\alpha_t = u_{1t} / 1.0000$
\end{enumerate}

\textbf{Hasil Akhir:}
\[ \boldsymbol{\alpha} = [0.4480, 0.4451, 0.4432, 0.4508, 0.4489] \]
\end{tcolorbox}

\begin{tcolorbox}[formula_box, title={16. Matriks Kompromi ($\mathbf{W}_c$)}]
\textbf{Formula:}
\[ \mathbf{W}_c = \sum_{t=1}^{T} \alpha_t \mathbf{W}_t \]

\textbf{Keterangan Variabel:}
\begin{itemize}
  \item $\mathbf{W}_c$: Rata-rata tertimbang dari semua matriks tahunan
  \item $\alpha_t$: Bobot optimal tahun ke-$t$
  \item $\text{tr}(\mathbf{W}_c) = $ total inersia dalam representasi kompromi
\end{itemize}

\textbf{Substitusi (Langkah demi Langkah):}
\begin{enumerate}
  \item $\mathbf{W}_c = 0.4480\mathbf{W}_{2020} + 0.4451\mathbf{W}_{2021} + 0.4432\mathbf{W}_{2022} + 0.4508\mathbf{W}_{2023} + 0.4489\mathbf{W}_{2024}$
  \item $\text{tr}(\mathbf{W}_c) = 91.6771$
  \item Nilai eigen pertama: $\lambda_1 = 45.0231$ $(49.1\%$ inersia)
  \item Nilai eigen kedua: $\lambda_2 = 14.7781$ $(16.1\%$ inersia)
\end{enumerate}

\textbf{Hasil Akhir:}
\[ \mathbf{W}_c \text{ berukuran } 41\times41, \quad \lambda_{1,2} = (45.0231,\ 14.7781) \]
\end{tcolorbox}

\begin{tcolorbox}[formula_box, title={17. Koordinat Skor Kompromi ($\mathbf{F}_c$)}]
\textbf{Formula:}
\[ \mathbf{F}_c = \mathbf{V}_c \mathbf{\Lambda}_c^{1/2} \]

\textbf{Keterangan Variabel:}
\begin{itemize}
  \item $\mathbf{V}_c$: Matriks eigenvector dari $\mathbf{W}_c$
  \item $\mathbf{\Lambda}_c$: Matriks diagonal eigenvalue
  \item Baris ke-$i$ dari $\mathbf{F}_c$ = koordinat biplot wilayah ke-$i$
\end{itemize}

\textbf{Substitusi (Langkah demi Langkah):}
\begin{enumerate}
  \item $\lambda_1 = 45.0231 \Rightarrow \sqrt{\lambda_1} = 6.7099$
  \item $\lambda_2 = 14.7781 \Rightarrow \sqrt{\lambda_2} = 3.8442$
  \item Kep.\ Seribu: $\mathbf{F}_c[1,:] = (-0.4150,\ 0.4857)$
  \item Kota Jak.\ Selatan: $\mathbf{F}_c[2,:] = (2.1105,\ 0.2253)$
\end{enumerate}

\textbf{Hasil Akhir:}
\[ \mathbf{F}_c \text{ berukuran } 41\times2 \text{ (dua PC dominan)} \]
\end{tcolorbox}

\begin{tcolorbox}[formula_box, title={29. Uji Permutasi Stabilitas}]
\textbf{Formula:}
\[p\text{-value}=\frac{\sum_{b=1}^B\mathbb{I}(\theta_1^{*(b)}\geq\theta_1^{obs})+1}{B+1}\]

\textbf{Keterangan Variabel:}
\begin{itemize}
  \item $\theta_1^{obs}$: Eigenvalue pertama $\mathbf{R}_{RV}$
  \item $B=999$ permutasi
  \item $p<0.05$: Struktur signifikan
\end{itemize}

\textbf{Substitusi (Langkah demi Langkah):}
\begin{enumerate}
  \item $\theta_1^{obs}=4.7825$
  \item 0/999 melebihi $\Rightarrow p=(0+1)/(999+1)$
\end{enumerate}

\textbf{Hasil Akhir:}
\[ p=0.001 \quad \text{(Tolak $H_0$: Stabil)} \]
\end{tcolorbox}


\section{Analisis Klaster dan Spasial}
\begin{tcolorbox}[formula_box, title={18. Jarak Euclidean antar Trajektori}]
\textbf{Formula:}
\[ d(i,j) = \sqrt{\sum_{t=1}^{T} \|\mathbf{f}_{i,t} - \mathbf{f}_{j,t}\|^2} \]

\textbf{Keterangan Variabel:}
\begin{itemize}
  \item $\mathbf{f}_{i,t}$: Koordinat PC wilayah $i$ pada tahun $t$ (dari proyeksi tahunan $\mathbf{F}_t = \mathbf{Z}_t \mathbf{V}_c$)
  \item $d(i,j)$: Jarak lintasan spasio-temporal antar dua wilayah
\end{itemize}

\textbf{Substitusi (Langkah demi Langkah):}
\begin{enumerate}
  \item Contoh: Kabupaten Bogor (i=6) vs Kota Bogor (i=25)
  \item Selisih koordinat per tahun dihitung, dikuadratkan, dijumlahkan
  \item $d = \sqrt{\sum_{t=2020}^{2024} \|\mathbf{f}_{Kab,t} - \mathbf{f}_{Kota,t}\|^2}$
\end{enumerate}

\textbf{Hasil Akhir:}
\[ d(\text{Kab.\ Bogor},\text{Kota Bogor}) = 2.1643 \]
\end{tcolorbox}

\begin{tcolorbox}[formula_box, title={19. Fungsi Objektif K-Means (SSE)}]
\textbf{Formula:}
\[ SSE = \sum_{j=1}^{k} \sum_{\mathbf{x}_i \in C_j} \|\mathbf{x}_i - \boldsymbol{\mu}_j\|^2 \]

\textbf{Keterangan Variabel:}
\begin{itemize}
  \item $k$: Jumlah klaster
  \item $C_j$: Himpunan anggota klaster ke-$j$
  \item $\boldsymbol{\mu}_j$: Centroid klaster ke-$j$ (rata-rata trajektori semua anggota)
  \item Minimasi SSE dilakukan secara iteratif oleh algoritma K-Means
\end{itemize}

\textbf{Substitusi (Langkah demi Langkah):}
\begin{enumerate}
  \item $k = 3$ klaster, berdasarkan trajektori $41$ wilayah selama $T=5$ tahun
  \item Klaster 1: 8 wilayah, Klaster 2: 13 wilayah, Klaster 3: 20 wilayah
  \item Setiap wilayah diwakili oleh vektor $\mathbf{x}_i$ berdimensi $2T = 10$
  \item $SSE = $ jumlah kuadrat jarak semua titik ke centroid klaster masing-masing
\end{enumerate}

\textbf{Hasil Akhir:}
\[ SSE = 25.0918 \quad \text{(setelah konvergensi K-Means)} \]
\end{tcolorbox}

\begin{tcolorbox}[formula_box, title={20. Statistik Global Moran's $I$}]
\textbf{Formula:}
\[ I = \frac{N}{\sum_i\sum_j w_{ij}} \cdot \frac{\sum_i\sum_j w_{ij}(x_i-\bar{x})(x_j-\bar{x})}{\sum_i(x_i-\bar{x})^2} \]

\textbf{Keterangan Variabel:}
\begin{itemize}
  \item $w_{ij}$: Bobot spasial (1 jika bertetangga, 0 jika tidak)
  \item $x_i$: Nilai PC1 wilayah ke-$i$
  \item $I > E[I]$: Autokorelasi spasial positif (klaster spasial)
  \item $E[I] = -1/(N-1)$: Ekspektasi di bawah hipotesis nol acak
\end{itemize}

\textbf{Substitusi (Langkah demi Langkah):}
\begin{enumerate}
  \item $N = 41$ wilayah (41 Kab/Kota di Banten, DKI Jakarta, Jawa Barat)
  \item Menggunakan skor PC1 kompromi sebagai variabel $x$
  \item Matriks bobot $\mathbf{W}$ berdasarkan keterkaitan geografis (Queen Contiguity)
  \item $E[I] = -1/(41-1) = -0.0250$
\end{enumerate}

\textbf{Hasil Akhir:}
\[ I_{obs} = 0.6903 \quad E[I] = -0.0250 \]
\textit{(Autokorelasi positif: wilayah dengan IPM tinggi cenderung bertetangga)}
\end{tcolorbox}

\begin{tcolorbox}[formula_box, title={26. Matriks Pembobot Spasial}]
\textbf{Formula:}
\[ w_{ij}=\begin{cases}1&\text{bertetangga}\\0&\text{selainnya}\end{cases},\; w_{ij}^*=\frac{w_{ij}}{\sum_k w_{ik}} \]

\textbf{Keterangan Variabel:}
\begin{itemize}
  \item $w_{ij}$: Bobot biner (Queen Contiguity)
  \item $w_{ij}^*$: Row-standardized
\end{itemize}

\textbf{Substitusi (Langkah demi Langkah):}
\begin{enumerate}
  \item $\mathbf{W}$: $41\times41$
  \item 5 tetangga $\Rightarrow w^*=0.2$
\end{enumerate}

\textbf{Hasil Akhir:}
\[ \mathbf{W} \text{ (41x41, simetri)} \]
\end{tcolorbox}

\begin{tcolorbox}[formula_box, title={27. Matriks Transisi Temporal}]
\textbf{Formula:}
\[\mathbf{T}=\begin{pmatrix}T_{11}&T_{12}&T_{13}\\T_{21}&T_{22}&T_{23}\\T_{31}&T_{32}&T_{33}\end{pmatrix}\]

\textbf{Keterangan Variabel:}
\begin{itemize}
  \item $T_{ab}$: Jumlah wilayah berpindah klaster $a\to b$
  \item Diagonal dominan $=$ stabil
\end{itemize}

\textbf{Substitusi (Langkah demi Langkah):}
\begin{enumerate}
  \item $k=3$, $41$ wilayah, $T=5$ tahun
\end{enumerate}

\textbf{Hasil Akhir:}
\[ \text{Matriks perpindahan klaster} \]
\end{tcolorbox}

\begin{tcolorbox}[formula_box, title={28. Ward's Linkage (Dendrogram)}]
\textbf{Formula:}
\[\Delta(A,B)=\frac{n_A n_B}{n_A+n_B}\|\boldsymbol{\mu}_A-\boldsymbol{\mu}_B\|^2\]

\textbf{Keterangan Variabel:}
\begin{itemize}
  \item $\Delta$: Peningkatan varians jika digabung
  \item $n_A,n_B$: Jumlah anggota
\end{itemize}

\textbf{Substitusi (Langkah demi Langkah):}
\begin{enumerate}
  \item $41$ wilayah, dipotong $k=3$
\end{enumerate}

\textbf{Hasil Akhir:}
\[ k=3 \text{ klaster} \]
\end{tcolorbox}

\end{document}
